\section*{Capítulo \ref{ch:b2:magmatic-series} --- Séries magmáticas e contextos tectônicos}

\begin{solution}{exo:b2:magmatic-series:1}
(a) $\ce{SiO2} = 49.60$, álcalis $2.22 + 0.51 = 2.73$: entre 45 e 52 e abaixo de 5,
\emph{basalto} (B). (b) 58.5 e $3.9 + 1.6 = 5.5$: entre 57 e 63; o limite
superior ali vai de (57, 5.9) a (63, 7.0), e está em $5.9 + 1.1\times1.5/6 = 6.18$
para 58.5; 5.5 fica abaixo dele: \emph{andesito} (O2). (c) 72.0 e $4.0 + 4.2 = 8.2$:
além da vertical em 69, acima de 8 e acima da reta de (69, 8) a (77, 0), que está
em 5.0 para 72: \emph{riolito} (R). (d) 47.0 e $4.5 + 2.5 = 7.0$: em 47, o campo U1
fica entre a reta de (45, 5) a (49.4, 7.3), em $5 + 2.3\times2/4.4 = 6.05$, e
a reta de (45, 9.4) a (49.4, 7.3), em $9.4 - 2.1\times2/4.4 = 8.45$. A rocha está
em U1: um \emph{tefrito} se sua norma contém menos de \qty{10}{\percent} de olivina, um
\emph{basanito} se contém mais.
\end{solution}

\begin{solution}{exo:b2:magmatic-series:2}
(a) Os óxidos somam 98.40 e o total com os voláteis, 100.00;
\ce{H2O} $1.40 < 2$ e \ce{CO2} $0.20 < 0.5$: a análise é aceitável. (b) Fator
$100/98.40 = 1.01626$: \ce{SiO2} 52.44, \ce{TiO2} 1.22, \ce{Al2O3} 16.36,
\ce{FeO_T} 9.96, \ce{MgO} 6.71, \ce{CaO} 8.33, \ce{Na2O} 3.56, \ce{K2O} 1.42
(soma 100.00); álcalis $3.56 + 1.42 = 4.98$. (c) Antes: 51.60 e 4.90, logo abaixo do
limite basalto--andesito basáltico em 52: \emph{basalto}. Depois: 52.44 e 4.98; em
52.44, o limite superior de O1, de (52, 5) a (57, 5.9), está em $5 + 0.9\times0.44/5
= 5.08$: \emph{andesito basáltico} (O1). O recálculo fez a rocha atravessar um
limite: os nomes perto de uma reta dependem dele. (d) Em 52.44, a reta de baixo K está em $0.3 +
0.045\times4.44 = 0.50$ e a de alto K em $1.2 + 0.095\times4.44 = 1.62$; 1.42
fica entre as duas: \emph{médio K}.
\end{solution}

\begin{solution}{exo:b2:magmatic-series:3}
(a) $Q + A + P = 90$: $Q = 25/90 = \qty{27.8}{\percent}$ e $100\,P/(A + P) =
100\times35/65 = 53.8$: campo 3b, um \emph{monzogranito}, um granito. A biotita e a
hornblenda não entram na soma QAPF. (b) $Q + A + P = 65$: $Q = 5/65 =
\qty{7.7}{\percent}$ (entre 5 e 20) e $100\times55/60 = 91.7$ (acima de 90): campo
10*, quartzo-diorito, quartzo-gabro ou quartzo-anortosito. O \omterm{def:b2:magmatic-series:qapf}{diagrama QAPF} sozinho não
escolhe entre eles; o plagioclásio e os minerais máficos (aqui
\qty{35}{\percent} de clinopiroxênio) o fazem.
\end{solution}

\begin{solution}{exo:b2:magmatic-series:4}
A reta de baixo K é $0.3 + 0.045(\ce{SiO2} - 48)$, a de alto K $1.2 +
0.095(\ce{SiO2} - 48)$. (a) Em 50: 0.39 e 1.39; $0.40$ fica logo acima da reta de baixo K:
\emph{médio K}. (b) Em 56: 0.66 e 1.96: \emph{médio K}. (c) Em 62: 0.93 e 2.53;
$2.90 > 2.53$: \emph{alto K}. (d) Em 66: 1.11 e 2.91; $1.00 < 1.11$: \emph{baixo K}.
\end{solution}

\begin{solution}{exo:b2:magmatic-series:5}
(a) Basalto: $A + F + M = 3.2 + 14.5 + 4.5 = 22.2$; $A = 14.4$, $F = 65.3$, $M = 20.3$.
Andesito: $5.05 + 6.1 + 3.8 = 14.95$; $A = 33.8$, $F = 40.8$, $M = 25.4$. (b) A linha
nunca vai além de $F = 54.5$; o basalto, com $F = 65.3$, fica do lado de F
qualquer que seja seu $A$: \emph{toleítico}, um ferrobasalto de uma tendência de Fenner. Para o
andesito, em $A = 33.8$ a linha fica entre (49.6, 35.0) e (52.5, 30.0), em $F =
49.6 + 2.9\times1.2/5 = 50.3$; $40.8 < 50.3$: \emph{cálcio-alcalino}. (c) A magnetita,
que retira o ferro férrico de um \omterm{def:b2:magma-emplacement:magma}{magma} oxidado desde cedo e impede $F/M$ de
subir tão depressa quanto em um toleíto.
\end{solution}

\begin{solution}{exo:b2:magmatic-series:6}
MORB: $\mathrm{Nb_N} = 4/0.658 = 6.08$, $\mathrm{La_N} = 4.12/0.648 = 6.36$, razão
0.956. \omterm{def:b2:elements-in-the-earth:morb-oib}{Basalto de ilha oceânica}: $38/0.658 = 57.8$, $30.8/0.648 = 47.5$, razão 1.215. Basalto
de arco: $3/0.658 = 4.56$, $6.1/0.648 = 9.41$, razão 0.484. Só o basalto de arco tem uma
anomalia de nióbio, uma depressão em Nb em um aranhograma: o nióbio é empobrecido em relação a
seu vizinho La: o \omterm{def:b2:magma-emplacement:magma}{magma} de arco recebeu elementos da placa, mas pouco
nióbio. O MORB é quase plano; o \omterm{def:b2:elements-in-the-earth:morb-oib}{basalto de ilha oceânica} é mais rico
em Nb do que em La.
\end{solution}

\begin{solution}{exo:b2:magmatic-series:7}
(a) $\mathrm{Ti} = 1.6\times0.5993\times10^4 = \qty{9589}{ppm}$; Ti/100 = 95.9, Zr
110, $3\mathrm Y = 96$; soma 301.9; coordenadas 31.8, 36.4, 31.8: campo B. (b) Ti
16\,780, logo 167.8, 220, 87; soma 474.8; 35.3, 46.3, 18.3: campo D. (c) Ti 4794, logo
47.9, 50, 60; soma 157.9; 30.3, 31.7, 38.0: campo B. O basalto (b) é um basalto intraplaca
(de uma ilha oceânica ou de um continente); (a) e (c), no campo B, podem ser MORB, toleítos
de arco ou basaltos cálcio-alcalinos: o diagrama não pode decidir, e é preciso outro (um
aranhograma, a anomalia de Nb).
\end{solution}

\begin{solution}{exo:b2:magmatic-series:8}
(a) MORB $142/164 = 0.866$, \omterm{def:b2:elements-in-the-earth:morb-oib}{basaltos de ilhas oceânicas} $4/173 = 0.023$, basaltos de arco $101/201 =
0.502$. (b) Com probabilidades a priori iguais, a regra de Bayes dá $0.866/(0.866 + 0.023 + 0.502) =
0.866/1.391 = 0.623$. (c) Que é muito improvável que o basalto seja um \omterm{def:b2:elements-in-the-earth:morb-oib}{basalto de ilha oceânica}
(probabilidade $0.023/1.391 = 0.017$), mas não se ele é um MORB ou um basalto de arco:
cerca de três chances em cinco contra duas.
\end{solution}

\begin{solution}{exo:b2:magmatic-series:9}
Mols por \qty{100}{g}: \ce{Na2O} $5.5/61.98 = 0.08874$, \ce{K2O} $4.5/94.20 =
0.04777$, \ce{Al2O3} $10.5/101.96 = 0.10298$. (a) $(0.08874 + 0.04777)/0.10298 =
0.13651/0.10298 = 1.326 > 1$: o riolito é \emph{peralcalino}. (b) O feldspato pode
receber $0.10298$ dos $0.13651$ mols de óxido alcalino; o excesso, $0.03353$, é
$0.03353/0.13651 = 0.246$, um quarto dos álcalis, que devem entrar em outros minerais.
\end{solution}

\begin{solution}{exo:b2:magmatic-series:10}
(a) $s = (\lambda_F - \lambda_M)/(\lambda_F - \lambda_A)$: para $\lambda_F = 0.5$, $(0.5
- 0.8)/(0.5 - 0.4) = -3.0$; para $0.7$, $(-0.1)/(0.3) = -0.333$. (b) Resolvendo
$s(\lambda_F - 0.4) = \lambda_F - 0.8$, obtém-se $\lambda_F = (0.4s - 0.8)/(s - 1)$:
Islândia, $s = -1.10$: $(-0.44 - 0.8)/(-2.10) = 0.590$; Cascatas, $s = -0.72$:
$(-0.288 - 0.8)/(-1.72) = 0.633$. (c) O \omterm{def:b2:magma-emplacement:magma}{magma} das Cascatas, com o $\lambda_F$
maior, perde seu ferro mais depressa: a magnetita o retira. (d) $\ln(F/M)$ cresce à
taxa $\lambda_M - \lambda_F$ e $\ln(A/F)$ à taxa $\lambda_F - \lambda_A$; assim, para
cada unidade ganha por $\ln(A/F)$, $\ln(F/M)$ ganha $-s$: 1.10 na Islândia, 0.72 nas
Cascatas. $F/M$ aumenta ao longo das duas tendências, mas mais depressa na toleítica, a tendência de
Fenner.
\end{solution}

\begin{solution}{exo:b2:magmatic-series:11}
(a) BHVO-2, mols por \qty{100}{g}: \ce{SiO2} 0.82557, \ce{Al2O3} 0.13182, \ce{FeO}
0.15521, \ce{MgO} 0.18007, \ce{CaO} 0.20328, \ce{Na2O} 0.03580, \ce{K2O} 0.00545.
Ortoclásio 0.00545 e albita 0.03580; \ce{Al2O3} restante $0.13182 - 0.04125 = 0.09057$,
todo na anortita com a mesma quantidade de \ce{CaO}; \ce{CaO} restante $0.11271$ no diopsídio;
$(\ce{Mg,Fe})\ce{O}$ restante $0.33528 - 0.11271 = 0.22257$ no hiperstênio. Sílica
necessária: $6\times0.00545 + 6\times0.03580 + 2\times0.09057 + 2\times0.11271 + 0.22257 =
0.87663$; disponível 0.82557; faltam 0.05106. Cada $(\ce{Mg,Fe})\ce{O}$ transferido do
hiperstênio para a olivina poupa meia sílica, no máximo $0.22257/2 = 0.11129$: transferir
$0.10212$ dele basta, e $0.12045$ fica no hiperstênio. A norma contém olivina
e hiperstênio, e nem quartzo nem nefelina: a rocha é saturada em sílica no
sentido da definição, um \emph{olivina-toleíto} com hiperstênio normativo, da
\omterm{def:b2:magmatic-series:series}{série toleítica}, como os que emite o estágio de escudo de um vulcão havaiano. (b) Mols: 0.75732, 0.14712, 0.16008, 0.21092, 0.17832, 0.06131,
0.01592. Or 0.01592, ab 0.06131, an 0.06988, di 0.10843, hy $0.37100 - 0.10843 =
0.26257$; sílica necessária $0.09552 + 0.36786 + 0.13976 + 0.21686 + 0.26257 = 1.08257$;
disponível 0.75732; faltam 0.32525. Todo o hiperstênio transformado em olivina poupa
0.13129, restando 0.19396; cada mol de \ce{Na2O} transferido da albita para a nefelina
libera 4 \ce{SiO2}, de modo que $0.19396/4 = 0.04849$
dos 0.06131 mols de \ce{Na2O} da albita se tornam nefelina. Nefelina na norma: a
rocha é subsaturada em sílica, um \emph{basalto alcalino}, da \omterm{def:b2:magmatic-series:series}{série alcalina}.
\end{solution}

\begin{solution}{exo:b2:magmatic-series:12}
(a) Basaltos de arco $116/141 = 0.823$, MORB $21/101 = 0.208$, \omterm{def:b2:elements-in-the-earth:morb-oib}{basaltos de ilhas oceânicas}
$3/137 = 0.022$. (b) $0.823/(0.823 + 0.208 + 0.022) = 0.823/1.053 = 0.782$.
(c) Numerador $0.1\times0.823 = 0.0823$; denominador $0.0823 + 0.6\times0.208 +
0.3\times0.022 = 0.0823 + 0.1248 + 0.0066 = 0.2137$; probabilidade $0.385$. (d) O
mesmo teste sobre a mesma rocha dá 0.78 ou 0.39 conforme a frequência de cada contexto
no lugar onde a rocha foi encontrada: um \omterm{def:b2:magmatic-series:discrimination}{diagrama de discriminação} dá uma verossimilhança, e o
contexto geológico deve fornecer o resto.
\end{solution}

\begin{solution}{pb:b2:magmatic-series:1}
\noindent\textbf{Parte I.}

\textbf{1.} Óxidos: lava 1 99.17, lava 2 98.65, lava 3 98.20; com os voláteis,
99.97, 99.80 e 99.45, todos perto de 100. A água (0.70, 1.10, 1.20) fica abaixo de
\qty{2}{\percent} e o dióxido de carbono (0.10, 0.05, 0.05) abaixo de \qty{0.5}{\percent}: as
três são aceitáveis.

\textbf{2.} Fator $100/99.17 = 1.00837$: \ce{SiO2} 50.82, \ce{Na2O} 2.92, \ce{K2O}
0.63, álcalis 3.55.

\textbf{3.} Lava 2: $100/98.65 = 1.01368$, \ce{SiO2} 59.40, \ce{Na2O} 3.95, \ce{K2O}
1.17, álcalis 5.12. Lava 3: $100/98.20 = 1.01833$, \ce{SiO2} 66.50, \ce{Na2O} 4.48,
\ce{K2O} 1.88, álcalis 6.36.

\textbf{4.} Lava 1 (50.82, 3.55): basalto. Lava 2 (59.40, 5.12; o limite superior de O2
ali está em $5.9 + 1.1\times2.40/6 = 6.34$): andesito. Lava 3 (66.50, 6.36; entre 63
e a reta de (69, 8) a (77, 0), abaixo da reta de (63, 7) a (69, 8)): dacito.

\textbf{5.} Brutas: (50.4, 3.52), (58.6, 5.05), (65.3, 6.25): os mesmos campos. O
recálculo desloca cada ponto de no máximo \qty{1.2}{\percent} de sílica e nada
muda aqui, ao contrário do \cref{exo:b2:magmatic-series:2}.

\textbf{6.} O \omterm{def:b2:magmatic-series:qapf}{diagrama QAPF}, a partir da moda medida. Uma rocha plutônica de composição
andesítica é rica em plagioclásio, com pouco feldspato alcalino ou quartzo: deve
cair no campo 10 ou 10*, um diorito ou um quartzo-diorito (sua moda decide).

\noindent\textbf{Parte II.}

\textbf{7.} Lava 1: retas em 0.43 e 1.47, $\ce{K2O} = 0.63$: médio K. Lava 2: 0.81
e 2.28, 1.17: médio K. Lava 3: 1.13 e 2.96, 1.88: médio K.

\textbf{8.} \ce{FeO_T} e \ce{MgO} livres de voláteis: 8.67 e 7.97; 6.18 e 3.85; 4.07
e 1.63. Lava 1: $3.55 + 8.67 + 7.97 = 20.19$; $A = 17.6$, $F = 42.9$, $M = 39.5$.
Lava 2: 15.15; 33.8, 40.8, 25.4. Lava 3: 12.06; 52.7, 33.7, 13.5.

\textbf{9.} Lava 1, $A = 17.6$: a linha entre (42.8, 17.1) e (46.9, 18.1) está em
$F = 42.8 + 4.1\times0.5/1.0 = 44.9$; $42.9 < 44.9$, lado cálcio-alcalino, perto da
linha. Lava 2, $A = 33.8$: $F = 49.6 + 2.9\times1.2/5 = 50.3$; $40.8 < 50.3$,
cálcio-alcalina. Lava 3, $A = 52.7$: $F = 40.8 - 2.9\times2.7/5 = 39.2$; $33.7 < 39.2$,
cálcio-alcalina.

\textbf{10.} $\ce{FeO_T}/\ce{MgO}$: $8.67/7.97 = 1.09$, $6.18/3.85 = 1.61$, $4.07/1.63 =
2.50$. A razão aumenta, como em todo basalto que se diferencia, mas o próprio ferro
cai de 8.67 para 4.07, e $F$ cai no triângulo de 42.9 para 33.7: nenhum
enriquecimento em ferro. Uma suíte toleítica subiria primeiro em direção a F, como faz o ferrobasalto do
\cref{exo:b2:magmatic-series:5}.

\textbf{11.} A \omterm{def:b2:magmatic-series:series}{série cálcio-alcalina}. A magnetita, cristalizando a partir de um \omterm{def:b2:magma-emplacement:magma}{magma} oxidado,
removeu o ferro junto com o magnésio; as espécies oxidadas vêm da placa que se desidrata em uma
zona de subducção.

\noindent\textbf{Parte III.}

\textbf{12.} $\mathrm{Nb_N} = 2.0/0.658 = 3.04$, $\mathrm{La_N} = 7/0.648 = 10.80$,
$\mathrm{Th_N} = 1.1/0.0795 = 13.84$.

\textbf{13.} $3.04/10.80 = 0.281$, bem abaixo das medianas dos MORB (0.956) e dos
\omterm{def:b2:elements-in-the-earth:morb-oib}{basaltos de ilhas oceânicas} (1.215), e abaixo até do basalto de arco mediano (0.484): uma profunda
depressão de nióbio.

\textbf{14.} $13.84/3.04 = 4.55$; o MORB mediano do aranhograma tem $2.64/6.08 =
0.43$. O tório, como os outros elementos enriquecidos nos basaltos de arco, foi acrescentado pela
placa em subducção.

\textbf{15.} \ce{TiO2} livre de voláteis $0.95\times1.00837 = 0.958$; $\mathrm{Ti} =
0.958\times5993 = \qty{5741}{ppm}$; Ti/100 = 57.4, Zr 70, $3\mathrm Y = 54$, soma 181.4:
31.6, 38.6, 29.8, no campo B, partilhado pelos MORB, toleítos de arco e basaltos cálcio-alcalinos.
Isso exclui um basalto intraplaca, mas não decide entre uma dorsal e um arco. O
diagrama foi construído para basaltos: em um andesito e em um dacito, o titânio foi removido
pelos óxidos e o zircônio concentrado pela cristalização, e seus campos (a lava 2
cairia em C, a lava 3 fora de todo campo) nada significariam.

\textbf{16.} Uma \omterm{def:b2:magmatic-series:series}{série cálcio-alcalina} de médio K, com uma profunda depressão de nióbio e
$\mathrm{Th_N}/\mathrm{Nb_N}$ alto, em um vulcão que emitiu de basalto a dacito: o vulcão
cresceu em um arco vulcânico acima de uma zona de subducção.

\noindent\textbf{Parte IV.}

\textbf{17.} $C_L = C_0/F$ para um elemento perfeitamente incompatível: $F = 0.63/1.17 =
0.538$.

\textbf{18.} Rubídio $10/20 = 0.500$; tório $1.1/2.1 = 0.524$. (Os \omterm{def:b2:elements-in-the-earth:trace}{elementos-traço},
dados tais como analisados, variariam menos de \qty{1}{\percent} com o
recálculo.)

\textbf{19.} $C_L/C_0 = F^{D - 1}$: $125/70 = 1.786 = F^{-0.9}$, $F = 1.786^{-1/0.9} =
0.525$.

\textbf{20.} Pela potassa, $F = 1.17/1.88 = 0.622$ da lava 2 à lava 3. Estrôncio:
$380/480 = 0.792 = 0.622^{D - 1}$, logo $D - 1 = \ln 0.792/\ln 0.622 = -0.234/-0.474 =
0.49$ e $D \approx 1.49$: o estrôncio se tornou compatível. O plagioclásio, que capta o
estrôncio, cristalizava em abundância.

\textbf{21.} Pela potassa, $F = 0.63/1.88 = 0.335$ (rubídio $10/32 = 0.31$, tório
$1.1/3.4 = 0.32$): cerca de dois terços do basalto, $1 - 0.335 = 0.665$, cristalizaram
antes do dacito.

\textbf{22.} A potassa assimilada elevaria o \ce{K2O} da lava 2 acima do que a
cristalização sozinha dá, de modo que $F = 0.63/1.17$ sairia baixo demais, e a massa de
cristais alta demais.

\textbf{23.} $(0.538 + 0.500 + 0.524 + 0.525)/4 = 2.087/4 = 0.522$; cristais $1 - 0.522
= 0.478$ do basalto, isto é, \textbf{\qty[text-series-to-math=true]{48}{kg} de
cristais por \qty{100}{kg} de basalto} para fazer o andesito.
\end{solution}
